\chapter{Reading a lit set on a boundary}

A set of lit points on a sphere admits three readings: where the points sit, how they push the boundary, and how they twist it. This chapter defines the three, gives the floor under each, and states what each can and cannot see. Nothing in the readings knows what a lit point means. The same readings serve a hash state, a solar system or a file of bytes without changing, and the caller supplies the meaning.

\section{Placements}

A placement decides where index \(k\) sits on the sphere. The choice shapes what can be read. Two are used and neither is a default.

\begin{center}
\begin{tabular}{@{}lp{0.3\linewidth}p{0.3\linewidth}@{}}
\toprule
placement & where index \(k\) goes & what it is good for \\
\midrule
golden, \(N\) points & height \(1 - 2(k+0.5)/N\), longitude \(k\gamma\)~\cite[Corollary~2]{src:Swierczkowski-1958} & even coverage with no seam and no pole pile~\cite{src:Gonzalez-2010} \\
rings, \(\mathit{rings} \times \mathit{width}\) & \(\mathit{rings}\) latitudes of equal width, \(\mathit{width}\) longitudes on each & a shift along a ring is a pure rotation \\
\bottomrule
\end{tabular}
\end{center}

The golden placement carries the property the screw below depends on: one step in the index is one fixed move, the same move everywhere along the spiral. The index shift \(n\) turns by \(n\gamma\) about the axis and slides the height by \(-2n/N\). That is rigid in height and longitude, and it keeps area by Archimedes' hat-box theorem~\cite[\emph{On the Sphere and Cylinder} I, Propositions 33, 42 and 43]{src:Heath-1897}. It is not an isometry of the sphere, because the axial radius \(\sqrt{1 - y^2}\) changes under the slide.

The ring placement puts its latitudes at the interiors of the bands. No ring lands on a pole where every point of it would pile into one place. The ring placement suits an operation that is a rotation, and the golden placement an operation that is a shift of the index.

\section{Readings}

The two harmonic readings come off one table, the boundary coefficients of the lit set summed once, and take that table apart in different directions. The third reading counts octants.

\begin{center}
\begin{tabular}{@{}lp{0.5\linewidth}@{}}
\toprule
reading & what it answers \\
\midrule
deflection & how much power the lit set carries at each degree \\
torsion & the phase of the same coefficients, by degree and order \\
octant share & what fraction of the set sits in each of the eight octants \\
\bottomrule
\end{tabular}
\end{center}

Moving every longitude by \(\alpha\) multiplies the entry at order \(m\) by \(e^{-i m \alpha}\). The magnitude holds still and the phase moves by \(m\alpha\). Deflection reads the magnitudes and torsion reads the phases, and a turn moves the phases while leaving the magnitudes alone.

\section{Choosing between deflection and torsion}

Deflection is blind to a rotation by construction~\cite{src:Kazhdan-Funkhouser-and-Rusinkiewicz-2003}. Power per degree is a magnitude, and a magnitude does not record how the object was turned. Deflection is invariant under every rotation, because a rotation acts on each degree by a unitary block~\cite[Section~2, eq.~(15); Section~3, eq.~(52)]{src:Gumerov-and-Duraiswami-2014}. Torsion measures the turn, and its sign gives the handedness; it recovers rotations about the polar axis, and those are the rotations tested.

Moving a lit set along its rings by seven different amounts and reading both:

\begin{center}
\begin{tabular}{@{}lp{0.5\linewidth}@{}}
\toprule
reading & what it does under a rotation \\
\midrule
deflection & moves by \(2.5 \times 10^{-14}\) of itself, which is nothing \\
torsion & recovers the angle to \(1.8 \times 10^{-13}\) radians, sign included \\
\bottomrule
\end{tabular}
\end{center}

The split is exact in both directions. The choice is not a matter of taste. A caller whose operations are rotations reads the twist. A caller wanting a quantity that holds still while the object turns reads the push. A caller who does not yet know which of those they have should read both and keep whichever carries a signal.

The recovered turn is in radians, in the direction the longitudes run. Order one pins the angle down without a wrap of its own; a higher order divides the angle and returns it only up to its own fraction of a turn. When neither set has a usable entry at the order asked, there is no turn to recover, a real outcome and not an error.

\section{The octant reading}

The octant share splits space into eight regions by the sign of each coordinate. Every region has a trilateral right-angle corner at the origin, all eight corners meet at that one point, and on the boundary the same split cuts eight congruent spherical triangles of area \(\pi/2\) each, by Girard~\cite[images f58 to f62]{src:Girard-1629}, totaling \(4\pi\). The measured shares sum to \(1.000000000000000\).

The eight sign frames are not all the same kind of move. Their determinants split four at \(+1\) and four at \(-1\). Half are rotations of the first octant and half carry a mirror. A caller reading handedness off this split should know it is there before relying on it.

The octant delta reports how much of the reading moved between two shares, as a percentage of the whole~\cite[Section~4.1, Proposition~4.2]{src:Levin-and-Peres-2017}. It is halved, because a share leaving one octant arrives in another and the two ends of a single move would otherwise be counted twice.

Eight regions is a small alphabet, and how small is a counting fact and not a judgment. The reading is rank \(8\), seven free numbers once the weight is fixed, and against a source with \(256\) degrees of freedom it moves in \(8\) directions and is blind in \(248\). A caller should know that number before building on this reading. Distinctness under such a reading may also be close to free: seven reals can separate a few dozen arbitrary states whether the seven carry anything. That last statement is a reading and is not derived here.

\section{The screw}

On a golden placement, shifting every index by a fixed amount comes to a rigid move: a turn, an axial slide, and the pitch, the slide per unit of turn. The pitch is \(-2/(N\gamma)\), and \(N = 256\) gives \(-0.003255258\) whatever the amount is. Computed across six amounts, the spread between the largest pitch and the smallest is \(4.3 \times 10^{-19}\), the rounding of the division; it says nothing about the placement. One axis and one pitch serve every shift, and a shift comes to a slide of the whole pattern along one fixed helix. The move is rigid in height and longitude and is not an isometry of the sphere.

The indices a shift runs off the end of return as a second rigid arm, a full axial extent away from where the screw alone would put them. A shift is the screw plus that arm, and where a lit point lands is fixed once the amount is known.

\section{What it costs}

Forming the coefficient table is the expensive step. It walks the lit set once per order and computes a Legendre column per point. The work grows as the lit count times the top degree squared. Degree eight over a few hundred lit points is immediate. A caller wanting a much higher degree should form the table once and hand it to both readings.

The Legendre column climbs the diagonal and then recurs in degree, keeping every intermediate at unit scale~\cite{src:Holmes-and-Featherstone-2002}. Computing the same values from factorials overflows and loses digits well before it overflows, and it fails quietly: the low degrees stay right while the fine structure goes wrong. The double's largest value is \(1.797 \times 10^{308}\). The normalization's \((l+m)!\) passes it at \(l + m = 171\), degree \(86\) at \(m = l\) and degree \(171\) at \(m = 0\). The unnormalized diagonal \((2l-1)!!\) passes it at \(l = 151\), since \(\log_{10} 299!! \approx 306.6\). The recurrence prevents a picture that looks plausible and is not.

\section{Choosing the degree}

The degree is a choice about how much a reading carries, and not a limit the arithmetic imposes. The basis measures orthonormal to \(2.2 \times 10^{-14}\) at degree fifteen under exact quadrature~\cite{src:Delsarte-Goethals-and-Seidel-1977}, and to \(1.8 \times 10^{-14}\) at degree eight. Precision is not what holds a reading at any degree a caller is likely to want. The overflow above is a real ceiling and it sits an order of magnitude beyond that.

What does bound the reading is counting. A reading to degree \(L\) carries exactly \((L+1)^2\) real numbers about its source: for a real function \(a_{l,-m} = (-1)^m \overline{a_{lm}}\), which leaves \((L+1)^2\) independent reals of the complex table's \(2(L+1)^2\). Deflection itself carries \(L+1\) numbers. Against a source with \(N\) degrees of freedom a reading to degree \(L\) is blind in \(N - (L+1)^2\) directions at any precision whatsoever.

\begin{center}
\begin{tabular}{@{}rrrrr@{}}
\toprule
degree & coefficients & rank against 256 & blind & least singular value \\
\midrule
8 & 81 & 81 & 175 & \(4.365\) \\
12 & 169 & 169 & 87 & \(3.818\) \\
15 & 256 & 256 & 0 & \(3.474 \times 10^{-3}\) \\
16 & 289 & 256 & 0 & \(1.524 \times 10^{-1}\) \\
\bottomrule
\end{tabular}
\end{center}

Pick the degree against the last column and never against the second. The smallest complete basis is degree fifteen, where \((L+1)^2\) first reaches \(256\), and it is the worst-conditioned usable one: the least singular value collapses to \(3.5 \times 10^{-3}\) there and recovers to \(1.5 \times 10^{-1}\) one degree higher, putting the worth of a single degree of headroom at \(44\) times in conditioning~\cite{src:Womersley-and-Sloan-2001}. A caller picking the cheapest complete basis pays for that saving in precision afterwards.

\section{The floors these readings sit on}

Every reading here has a move that cannot change it, and running that move gives the reading's own grain in the units it reports. A threshold is then measured and never picked.

\begin{center}
\begin{tabular}{@{}p{0.22\linewidth}p{0.42\linewidth}r@{}}
\toprule
reading & the move that cannot change it & residual \\
\midrule
deflection & a rotation of the whole set & \(2.539 \times 10^{-14}\) \\
torsion & a rotation by a whole turn & exactly 0 \\
octant share & a lit point relabeled inside its own octant & exactly 0 \\
octant delta & a reading against itself & exactly 0 \\
screw pitch & changing the shift amount & \(4.337 \times 10^{-19}\) \\
spectrum power & a rotation of the sources & \(4.005 \times 10^{-16}\) \\
a reading over arms & the arms redrawn as different shapes at the same weight & \(2.220 \times 10^{-16}\) \\
\bottomrule
\end{tabular}
\end{center}

The last row is a null over the reading instead of over the object. An arm is identified by its topology together with its weight, and a shape is one realization of that. Drawing the arms differently while holding the weight cannot move a letter. One arm set drawn four ways, including as a line with a dwell particle on it, leaves every placement point on its arm.

Three of the seven are exactly zero and not merely small, because the moves are exact on integers. The octant reading has no continuum of small changes at all: a share is a count over a weight and both are integers, and at fixed weight the smallest change it can register is \(1/\mathit{weight}\). Two states of \emph{differing} weight can still give shares arbitrarily close together, leaving a caller who compares across weights to make a modeling choice about which states count as the same. That choice is a modeling decision and not a tolerance.

The companion measurement is not a null. Torsion recovers a known angle to \(1.803 \times 10^{-13}\) radians, which asks whether the reading moves the right amount when something did happen. A null asks whether it holds still when nothing did. A reading wants both answers and they are different questions.

\section{Testing the floors}

Before it reports any floor, the floor measurement hands itself a move that is not a null while declaring that it is, and a genuine null of the same shape. It has to catch the first and stay quiet on the second, and it withholds every floor if either misbehaves. A checker that has never caught its own defect is untested. The redraw floor is measured at three levels, and four drawings that are not redraws are held against it first.

The readings, the placements and the floor measurement are implemented in Python with the standard library alone, and are in the orior repository, \texttt{https://github.com/dstroy0/orior}.
